\section{2015级工科数学分析下 A卷}

\subsection{资料来源}
\begin{itemize}
  \item 试题：2015级软件工科数学分析下 A 卷，DOC 正文已抽取；公式对象已按 Word 公式图片逐项恢复。
  \item 公式图片审计：\texttt{tmp/past-exam-equations/2015-A/manifest.csv}。
\end{itemize}

\subsection{试题}
诚信应考，考试作弊将带来严重后果！

华南理工大学本科生期末考试

《工科数学分析》2015--2016 学年第二学期期末考试试卷（A）卷

注意事项：
\begin{enumerate}
  \item 开考前请将密封线内各项信息填写清楚；
  \item 所有答案请直接答在试卷上（或答题纸上）；
  \item 考试形式：闭卷；
  \item 本试卷共 5 个大题，满分 100 分，考试时间 120 分钟。
\end{enumerate}

\subsubsection*{一、填空题（每小题 3 分，共 15 分）}
\begin{enumerate}
  \item 设 $y_1=3,\ y_2=3+x^2,\ y_3=3+x^2+e^x$ 都是方程
  \[
  (x^2-2x)y''-(x^2-2)y'+(2x-2)y=6x-6
  \]
  的解，则该方程的通解为 \underline{\hspace{4cm}}。
  \item 设由方程 $F(x-y,y-z,z-x)=0$ 确定了 $z=z(x,y)$，其中 $F$ 可微，则 $dz=$ \underline{\hspace{3cm}}。
  \item 函数
  \[
    f(x,y,z)=\ln(x^2+y^2+z^2)
  \]
  在点 $P(2,0,1)$ 处沿梯度方向的方向导数是 \underline{\hspace{4cm}}。
  \item 设 $L:\dfrac{x^2}{a^2}+\dfrac{y^2}{b^2}=1$ 取顺时针方向，则
  \[
  \int_L(x-y)\,dx+(x+y)\,dy=
  \]
  \underline{\hspace{3cm}}。
  \item 设 $f(x)=\sin(\pi x^2)\ (0\le x<1)$，而
  \[
  S(x)=\sum_{n=1}^{\infty}b_n\sin(n\pi x)\quad(-\infty<x<+\infty),
  \]
  其中 $b_n=2\int_0^1 f(x)\sin(n\pi x)\,dx\ (n=1,2,\cdots)$，则 $S\!\left(-\dfrac12\right)$ \underline{\hspace{3cm}}。
\end{enumerate}

\subsubsection*{二、计算题（每小题 8 分，共 40 分）}
\begin{enumerate}
  \item 设函数 $z=f(x^2-y^2,2xy)$，其中 $f(\xi,\eta)$ 具有连续的二阶偏导数，且满足
  \[
  \frac{\partial^2 f}{\partial \xi^2}+\frac{\partial^2 f}{\partial \eta^2}=0,
  \]
  求 $\dfrac{\partial^2 z}{\partial x^2}+\dfrac{\partial^2 z}{\partial y^2}$。
  \item 计算曲线积分
  \[
  I=\int_L\sqrt{x^2+y^2}\,ds,
  \]
  其中 $L$ 是圆周 $x^2+y^2=ax$。
  \item 设 $f(x)$ 在 $(-\infty,+\infty)$ 内有连续导数，$L$ 是从点 $A\!\left(3,\dfrac23\right)$ 到点 $B(1,2)$ 的直线段，计算曲线积分
  \[
  I=\int_L\frac{1+y^2f(x)}{y}\,dx+\frac{x}{y^2}\bigl[y^2f(x)-1\bigr]\,dy.
  \]
  \item 设 $\Sigma$ 为曲面 $x^2+y^2+z^2=2ax\ (a>0)$，计算曲面积分
  \[
  I=\iint_\Sigma (x^2+y^2+z^2)\,dS.
  \]
  \item 设 $\Sigma$ 是上半球面 $z=\sqrt{a^2-x^2-y^2}\ (a>0)$，取上侧，计算曲面积分
  \[
  I=\iint_\Sigma x^3\,dy\,dz+y^3\,dz\,dx+z^3\,dx\,dy.
  \]
\end{enumerate}

\subsubsection*{三、解答题（每小题 9 分，共 18 分）}
\begin{enumerate}
  \item 求幂级数
  \[
  \sum_{n=1}^{\infty}\frac{2n+1}{n!}x^{2n}
  \]
  的和函数。
  \item 求微分方程
  \[
  xy'-2y=x^3\sin x
  \]
  的通解。
\end{enumerate}

\subsubsection*{四、证明题（每小题 9 分，共 18 分）}
\begin{enumerate}
  \item 设 $f(x)\in C[0,1]$，求证：
  \[
  \int_0^1 dx\int_x^1 f(x)f(y)\,dy=\frac12\left[\int_0^1 f(x)\,dx\right]^2.
  \]
  \item 证明：函数项级数
  \[
  \sum_{n=1}^{\infty}\frac{\cos(nx)}{n^2}
  \]
  在 $(-\infty,+\infty)$ 上一致收敛。
\end{enumerate}

\subsubsection*{五、应用题（本题 9 分）}
求内接于半径为 $a$ 的球且有最大体积的长方体。

\subsection{答案与解析}

\subsubsection*{一、填空题}
\begin{enumerate}
  \item 因 $y_2-y_1=x^2,\ y_3-y_2=e^x$ 是对应齐次方程的两个线性无关解，且 $y_1=3$ 是一个特解，故通解为
  \[
  y=3+C_1x^2+C_2e^x.
  \]
  \item 记 $F_1,F_2,F_3$ 分别为 $F$ 对三个自变量的偏导数。由
  \[
  F_1(dx-dy)+F_2(dy-dz)+F_3(dz-dx)=0
  \]
  得
  \[
  dz=\frac{(F_3-F_1)\,dx+(F_1-F_2)\,dy}{F_3-F_2}.
  \]
  \item $\nabla f=\dfrac{2(x,y,z)}{x^2+y^2+z^2}$，故
  \[
  |\nabla f(2,0,1)|=\frac{2\sqrt5}{5}.
  \]
  \item 由 Green 公式，逆时针时积分为
  \[
  \iint_D\left(\frac{\partial}{\partial x}(x+y)-\frac{\partial}{\partial y}(x-y)\right)\,dA
  =2\pi ab.
  \]
  本题 $L$ 取顺时针方向，故结果为 $-2\pi ab$。
  \item 该正弦级数对应 $[0,1]$ 上的奇延拓且周期为 $2$，所以
  \[
  S\!\left(-\frac12\right)=-f\!\left(\frac12\right)=-\sin\frac{\pi}{4}=-\frac{\sqrt2}{2}.
  \]
\end{enumerate}

\subsubsection*{二、计算题}
\begin{enumerate}
  \item 令 $\xi=x^2-y^2,\ \eta=2xy$。该变换满足
  \[
  \xi_x^2+\eta_x^2=\xi_y^2+\eta_y^2=4(x^2+y^2),\qquad
  \xi_x\xi_y+\eta_x\eta_y=0,
  \]
  且 $\xi_{xx}+\xi_{yy}=\eta_{xx}+\eta_{yy}=0$。因此
  \[
  z_{xx}+z_{yy}=4(x^2+y^2)(f_{\xi\xi}+f_{\eta\eta})=0.
  \]
  \item 圆 $x^2+y^2=ax$ 可参数化为
  \[
  x=\frac a2(1+\cos t),\qquad y=\frac a2\sin t,\qquad 0\le t\le 2\pi.
  \]
  此时 $\sqrt{x^2+y^2}=a|\cos(t/2)|,\ ds=\dfrac a2\,dt$，故
  \[
  I=\frac{a^2}{2}\int_0^{2\pi}\left|\cos\frac t2\right|dt=2a^2.
  \]
  \item 直线段 $AB$ 上
  \[
  y=\frac{8-2x}{3},\qquad x:3\to 1.
  \]
  被积表达式可写为
  \[
  d\!\left(\frac{x}{y}\right)+f(x)\,d(xy).
  \]
  于是
  \[
  I=\left.\frac{x}{y}\right|_A^B+\int_3^1 f(x)\frac{8-4x}{3}\,dx
  =-4+\frac13\int_3^1(8-4x)f(x)\,dx.
  \]
  \item 曲面是以 $(a,0,0)$ 为球心、半径为 $a$ 的球面。在该球面上
  \[
  x^2+y^2+z^2=2ax.
  \]
  由球面对称性 $\iint_\Sigma x\,dS=a\cdot 4\pi a^2$，故
  \[
  I=2a\iint_\Sigma x\,dS=8\pi a^4.
  \]
  \item 把上半球面与底面圆盘合成上半球体的闭曲面。底面上 $z=0$，第三项通量为 $0$。由 Gauss 公式，
  \[
  I=\iiint_\Omega 3(x^2+y^2+z^2)\,dV
  =3\int_0^a r^2\cdot 2\pi r^2\,dr
  =\frac{6\pi a^5}{5}.
  \]
\end{enumerate}

\subsubsection*{三、解答题}
\begin{enumerate}
  \item 令 $t=x^2$，则
  \[
  \sum_{n=1}^{\infty}\frac{2n+1}{n!}x^{2n}
  =2\sum_{n=1}^{\infty}\frac{n t^n}{n!}+\sum_{n=1}^{\infty}\frac{t^n}{n!}
  =2te^t+(e^t-1).
  \]
  因此和函数为
  \[
  (2x^2+1)e^{x^2}-1.
  \]
  \item 当 $x\ne0$ 时，方程化为
  \[
  y'-\frac2x y=x^2\sin x.
  \]
  取积分因子 $x^{-2}$，得
  \[
  \left(\frac{y}{x^2}\right)'=\sin x.
  \]
  因而
  \[
  \frac{y}{x^2}=-\cos x+C,\qquad y=Cx^2-x^2\cos x.
  \]
\end{enumerate}

\subsubsection*{四、证明题}
\begin{enumerate}
  \item 记 $I=\int_0^1 dx\int_x^1 f(x)f(y)\,dy$。交换 $x,y$ 的角色可得
  \[
  I=\int_0^1 dy\int_0^y f(x)f(y)\,dx.
  \]
  两式相加覆盖单位正方形除对角线外的两个三角区域，而对角线面积为零，所以
  \[
  2I=\int_0^1\int_0^1 f(x)f(y)\,dx\,dy
  =\left[\int_0^1 f(x)\,dx\right]^2.
  \]
  结论得证。
  \item 对任意 $x\in(-\infty,+\infty)$，有
  \[
  \left|\frac{\cos(nx)}{n^2}\right|\le \frac1{n^2}.
  \]
  又 $\sum_{n=1}^{\infty}\dfrac1{n^2}$ 收敛，由 Weierstrass 判别法，原函数项级数在 $(-\infty,+\infty)$ 上一致收敛。
\end{enumerate}

\subsubsection*{五、应用题}
设长方体中心与球心重合，三个半棱长为 $x,y,z>0$。约束为
\[
x^2+y^2+z^2=a^2,\qquad V=8xyz.
\]
由对称性或 Lagrange 乘子法，最大值在 $x=y=z=\dfrac a{\sqrt3}$ 处取得。因此最大体积长方体为正方体，棱长为
\[
\frac{2a}{\sqrt3},
\]
最大体积为
\[
V_{\max}=\left(\frac{2a}{\sqrt3}\right)^3=\frac{8a^3}{3\sqrt3}.
\]
